\section{One number for the whole smile}
\label{sec:vsnumber}

Chapters~2--4 built three ways of fitting one expiry's smile, and each
chapter ended the same way: with a curve that reprices every quote to a
few basis points.  This section introduces a single number, attached to
the whole smile at once, that will spend the rest of the chapter
testing those curves in the region the quotes never reach.

The number comes from a real instrument.  A \emph{variance swap} is a
contract on how much the underlying actually moves.  Over the life of
the contract, each day's log-return is squared, the squares are summed,
and at expiry the buyer receives that realized sum and pays a fixed
amount agreed at the start.  The fixed amount is called the
\emph{strike} of the swap, and the strike at which the contract is
worth zero today is the \emph{fair strike}.  Quoting desks publish fair
strikes the way they publish option prices.  So the fair strike is not
a modeling construct: it is a market number, one per (underlier,
expiry), and a fitted smile implies a value for it.

To compute that implied value, work in the book's standing units.  The
underlying is measured as $Y=S/F$, the level divided by the forward, so
$Y$ starts at $1$ and is a positive martingale --- a process with no
drift, whose expected future value is its current value
(\cref{sec:lvfield}).  Time is variance time $\tau$ --- the maturity
variable the book has quoted variance against since Chapter~2, whose
clock Chapter~8 will make precise --- so a smile's total implied
variance at log-moneyness $k$ is $w(k)=\sigma(k)^2\tau$, with
$\sigma(k)$ the implied volatility there.  Let the martingale move as
a diffusion,
\[
\dd Y_s=\sqrt{v_s}\;Y_s\,\dd W_s,
\]
where $s$ runs from $0$ to the expiry $\tau$ and $v_s$ is the
\emph{instantaneous variance} at time $s$ --- the squared volatility
the process is running at that instant.  In the representative model of
Chapter~4, $v_s$ is the field evaluated along the path,
$v_s=v(s,Y_s)$; but nothing below needs that: $v_s$ may be any
reasonable random process.  As the monitoring interval shrinks, the
sum of squared returns converges to the accumulated instantaneous
variance, so in these units the realized variance the swap pays is
the integral $\int_0^{\tau}v_s\,\dd s$, and the fair strike is its
expected value under the pricing law --- the law the smile encodes,
the same one behind every expectation in this book.  (Two conventions
separate this from the listed contract, and it is worth stating them
once, here.  The contract sums squared \emph{daily} returns while the
integral monitors continuously, and the contract's clock is the
calendar while ours is variance time; both refinements have standard
corrections and Chapter~8 returns to the clock.  A third gap ---
paths that jump --- is genuine and reappears in the chapter's closing
table of guarantees, \cref{sec:vscontract}.)

The expected value looks like it should depend on the whole path law
of $Y$.  It does not, and the two-line computation that shows this is
the heart of the chapter.  Apply It\^o's formula --- the chain rule
for diffusions, which adds a second-derivative term --- to $\log Y_s$,
using the multiplication rule $(\dd W_s)^2=\dd s$, so that
$(\dd Y_s)^2=v_s Y_s^2\,\dd s$:
\begin{equation}
\dd \log Y_s=\frac{\dd Y_s}{Y_s}
-\frac12\,\frac{(\dd Y_s)^2}{Y_s^2}
=\sqrt{v_s}\,\dd W_s-\frac12\,v_s\,\dd s .
\label{eq:vsito}
\end{equation}
Integrate from $0$ to $\tau$ and use $\log Y_0=0$:
\[
\log Y_\tau=\int_0^{\tau}\sqrt{v_s}\,\dd W_s
-\frac12\int_0^{\tau}v_s\,\dd s .
\]
Now take expectations.  The first integral is a stochastic integral
against a Brownian motion, and such integrals have expectation zero ---
they are the running gains of a bet with no edge.  Writing
$X=\log Y_\tau$ for the terminal log-return, the surviving terms
rearrange to
\begin{equation}
w_{\rm vs}
\;=\;\E\!\left[\int_0^{\tau}v_s\,\dd s\right]
\;=\;-2\,\E[X].
\label{eq:vslogcontract}
\end{equation}
The expected realized variance equals minus twice the expected
log-return.  The right-hand side no longer mentions the path: it is a
property of the \emph{terminal law alone}, the same law every smile of
this book encodes.  A path-dependent payoff has been priced by a
European one --- the payoff $-2\log y$, called the \emph{log
contract}.  Chapter~2 already met this number in passing, as
\eqref{eq:varswap}; this chapter is about taking it seriously.  We
write $w_{\rm vs}$ for the fair strike in total-variance units and
$\sigma_{\rm vs}=\sqrt{w_{\rm vs}/\tau}$ for the same number quoted as
a volatility.

Before asking what the number says about a real smile, compute it once
by hand on the simplest smile there is.  Suppose the smile is flat:
$w(k)=w_0$ at every strike.  A flat smile means the law of $X$ is
exactly the one Black--Scholes assumes, a normal with variance $w_0$.
Its mean is pinned by the martingale condition
$\E[e^X]=\E[Y_\tau]=Y_0=1$: for a normal with mean $\mu_X$ and
variance $w_0$, the standard identity
$\E[e^X]=e^{\mu_X+w_0/2}$ forces
\[
\mu_X+\frac{w_0}{2}=0
\qquad\Longrightarrow\qquad
\E[X]=-\frac{w_0}{2},
\]
and \eqref{eq:vslogcontract} gives
\begin{equation}
w_{\rm vs}=-2\,\E[X]=w_0 .
\label{eq:vsflat}
\end{equation}
On a flat smile the fair strike \emph{equals} the ATM total variance,
exactly.  This small computation earns its keep twice over.  It is a
sanity check --- a variance swap on a constant-volatility world is
worth that variance.  And it is a decomposition: whatever a real fair
strike shows \emph{beyond} the ATM level is entirely a statement about
the smile's departure from flatness --- its curvature and, above all,
its wings.  \Cref{sec:vsterm} will measure that excess on real data
and find it worth hundreds of basis points.

Why the wings, specifically?  Look at the shape of the payoff.  The
function $-2\log y$ is large where $y$ is far from $1$, and it grows
without bound as $y\to0$.  At the same time the martingale condition
$\E[e^X]=1$ forces a balance: a law cannot put weight on large upside
without compensating weight below, and every rearrangement of tail
mass moves $\E[X]$.  The fair strike therefore loads exactly on the
part of the law the quoted strikes describe least --- a suspicion the
next two sections make precise and then quantify.

\Cref{fig:vspins} shows why the number deserves a chapter.  It fits
the three families of Part~I --- LQD from Chapter~2, SVI and MCS from
Chapter~3 --- to the same \MacVsPinsNQuotes\ mid quotes (bid--ask
midpoints) of the book's running node: the (underlier, expiry) pair
SPY December 2026 that Chapters~2--4 kept returning to.  The fits
follow Chapter~3's comparison protocol (\cref{sec:comparison}) ---
same quotes, equal weights, each family's reference settings.  In
panel~(a), the quoted span: the three fitted curves pass through the
same quotes and through each other --- the worst fit error of any
family is \MacVsPinsWorstRmsBp\,bp, and no two of the curves are ever
farther than \MacVsPinsBellyAgreeBp\,bp apart on the span (one vol bp
is $10^{-4}$ of absolute volatility, the chapter's unit for
volatility distances).  By every measure the data offers, the three models
are the same smile.  In panel~(b), the same three fits drawn wide, in
total variance, with the quoted span shaded: outside the shading the
curves separate --- each family continues the smile according to its
own structure --- and the legend prints each family's implied fair
strike.  The three values of $\sigma_{\rm vs}$ span
\MacVsPinsMaxGapBp\ vol bp: LQD says \MacVsPinsLqdPct\%, SVI says
\MacVsPinsSviPct\%, MCS says \MacVsPinsMcsPct\%.  Nothing has gone
wrong.  Three valid laws thread the same quotes and integrate to
different numbers.

\begin{figure}[htbp]
\centering
\paperfig{\textwidth}{fig_vs_pins.pdf}
\caption{Three families, one quote set, three fair strikes (SPY
December 2026, \MacVsPinsNQuotes\ quotes).  (a)~On the quoted span the
LQD, SVI and MCS fits agree with the quotes and with each other to
\MacVsPinsBellyAgreeBp\,bp.  (b)~In total variance, wide: beyond the
shaded quoted span the three curves separate, and their fair var-swap
strikes differ by up to \MacVsPinsMaxGapBp\ vol bp.}
\label{fig:vspins}
\end{figure}

We now have the chapter's protagonist and its puzzle.  The fair strike
is one market-traded number that the whole smile determines --- and
that fitted smiles, in perfect agreement wherever quotes exist, do not
agree on.  The plan follows the thesis of the book.
\Cref{sec:vsthree} computes the number exactly, in each model's own
coordinates, and checks that all routes agree on one model at a time.
\Cref{sec:vslives} locates the disagreement: it lives precisely in the
share of the integral that accrues beyond the last quote.
\Cref{sec:vsceiling,sec:vsenvelope} then prove what \emph{can} be said
with no quotes at all --- a ceiling and an envelope --- and
\cref{sec:vswingchoice,sec:vsterm} treat what remains inside those
bounds for what it is: a convention to be stated, policed at finite
strikes, and kept consistent across the term structure.
